\documentclass[11pt,a4paper]{article}

% ─── Packages ───
\usepackage[utf8]{inputenc}
\usepackage[T1]{fontenc}
\usepackage{amsmath,amssymb,amsthm}
\usepackage{booktabs}
\usepackage{array}
\usepackage{longtable}
\usepackage{graphicx}
\usepackage[colorlinks=true,linkcolor=blue,citecolor=blue,urlcolor=blue]{hyperref}
\usepackage{geometry}
\geometry{margin=1in}
\usepackage{enumitem}
\usepackage{listings}
\usepackage{xcolor}
\usepackage{natbib}
\bibliographystyle{plainnat}

% ─── Code listing style ───
\lstset{
  language=JavaScript,
  basicstyle=\ttfamily\small,
  keywordstyle=\color{blue!70!black},
  commentstyle=\color{green!50!black},
  stringstyle=\color{red!60!black},
  numbers=left,
  numberstyle=\tiny\color{gray},
  breaklines=true,
  frame=single,
  backgroundcolor=\color{gray!5},
}

% ─── Theorems ───
\newtheorem{theorem}{Theorem}
\newtheorem{proposition}[theorem]{Proposition}
\newtheorem{remark}{Remark}
\newtheorem{definition}{Definition}

% ─── Title ───
\title{%
  The Cycle Structure of the King Wen Permutation:\\
  A Group-Theoretic Analysis of Two Classical Hexagram Orderings%
}
\author{
  Zhengwen Ge\\
  \texttt{a15067006969@gmail.com}
}
\date{March 2026}

\begin{document}
\maketitle

% ═══════════════════════════════════════════
\begin{abstract}
The 64 hexagrams of the \textit{I Ching} (Book of Changes) admit two canonical orderings: the \emph{binary natural order}, a 6-bit counting sequence reconstructed by Shao Yong (1011--1077~CE) and retroactively attributed to the mythical Fu Xi; and the \emph{King Wen sequence}, the textual ordering of the received \textit{Zhouyi} ($\sim$1000~BCE).
Treating the map from the binary natural order to the King Wen sequence as a permutation $\sigma \in S_{64}$, we compute its cycle decomposition and obtain the cycle type $\mathbf{(52,\,10,\,2)}$ with zero fixed points.
This means 81.25\% of hexagrams are locked in a single cycle, no hexagram remains in its natural binary position, and the permutation order is $\mathrm{lcm}(52,10,2) = 260$.
We provide statistical context (comparison with random permutations of 64 elements), a Hamming distance analysis of the King Wen sequence, and a discussion of the epistemological boundaries of interpreting cycle membership semantically.
To our knowledge, this cycle decomposition has not been reported in the existing literature on \textit{I Ching} mathematics.
All results are computationally reproducible from the code provided in Appendix~A.
\end{abstract}

\bigskip
\noindent\textbf{Keywords:} permutation group, cycle decomposition, \textit{I Ching}, King Wen sequence, hexagram ordering, combinatorics

\noindent\textbf{MSC 2020:} 05A05, 20B30, 01A25


% ═══════════════════════════════════════════
\section{Introduction}\label{sec:intro}

The sixty-four hexagrams (\textit{guà} 卦) are the foundational symbols of the \textit{I Ching}, one of the oldest extant Chinese texts.
Each hexagram is a stack of six binary lines---solid (\textit{yáng} 阳, encoded as~1) or broken (\textit{yīn} 阴, encoded as~0)---yielding $2^6 = 64$ possible figures.

Two orderings of these 64 hexagrams have canonical status:

\begin{enumerate}[label=(\roman*)]
  \item \textbf{The binary natural order.}
  Reading each hexagram as a 6-bit integer from bottom to top, this is simply the sequence $0, 1, 2, \ldots, 63$.
  It was formalized by the Song-dynasty scholar Shao Yong (邵雍, 1011--1077~CE) in his \textit{Huángjí Jīngshì} (皇极经世) and retroactively attributed to the legendary sage Fu Xi.
  Leibniz recognized its binary nature in 1703 \citep{leibniz1703}.

  \item \textbf{The King Wen sequence.}
  This is the ordering found in the received text of the \textit{Zhouyi} (周易), traditionally attributed to King Wen of Zhou ($\sim$1100~BCE).
  Its organizational logic has been debated for millennia---Han Kangbo called it ``not the essence of the \textit{Yi},'' Zhu Xi deemed it ``unknowable'' (paraphrased from traditional commentary), and Gao Heng argued it is meaningful---but no consensus exists \citep{rutt2002,shaughnessy2014}.
\end{enumerate}

\begin{remark}[Historical clarification]\label{rem:historical}
  The binary natural order was formalized approximately 2{,}100 years \emph{after} King Wen.
  We do not claim or assume that King Wen ``rearranged'' the binary natural order.
  Our analysis treats the two orderings as independently defined mathematical objects and characterizes the permutation between them, analogous to comparing two genome sequences without presupposing which is ancestral.
\end{remark}

Despite extensive scholarly attention to the King Wen sequence, few studies have treated the relationship between these two orderings as a problem in pure algebra.
The present paper performs this elementary but apparently overlooked operation: we view the map from the binary natural order to the King Wen sequence as a permutation $\sigma \in S_{64}$ and compute its \emph{cycle decomposition}.


% ═══════════════════════════════════════════
\section{Definitions and Data}\label{sec:data}

\subsection{Binary Encoding}

We encode each hexagram as a 6-bit string read from the bottom line (least significant bit) upward.
The hexagram $\texttt{000000}$ (all broken lines) is Hexagram~2 (\textit{Kūn} 坤) in the King Wen numbering, and $\texttt{111111}$ (all solid lines) is Hexagram~1 (\textit{Qián} 乾).
The binary natural order is the sequence of hexagrams indexed $0$ through~$63$ by their integer value.

\subsection{The Permutation $\sigma$}

Let $W(i)$ denote the King Wen number (1-indexed) of the hexagram at binary natural position~$i$ (0-indexed).
We define the permutation $\sigma\colon \{0,1,\ldots,63\} \to \{0,1,\ldots,63\}$ by
\begin{equation}\label{eq:sigma}
  \sigma(i) = W(i) - 1.
\end{equation}
This maps each hexagram from its position in the binary natural order to its position in the King Wen order (both 0-indexed).
The complete mapping is given in Table~\ref{tab:mapping}.

\begin{table}[ht]
\centering
\caption{Binary natural position $i$ $\to$ King Wen number $W(i)$ for all 64 hexagrams.}
\label{tab:mapping}
\small
\begin{tabular}{@{}cccc|cccc@{}}
\toprule
$i$ & Binary & $W(i)$ & Name & $i$ & Binary & $W(i)$ & Name \\
\midrule
0 & 000000 & 2 & Kūn & 32 & 100000 & 24 & Fù \\
1 & 000001 & 23 & Bō & 33 & 100001 & 27 & Yí \\
2 & 000010 & 8 & Bǐ & 34 & 100010 & 3 & Zhūn \\
3 & 000011 & 20 & Guān & 35 & 100011 & 42 & Yì \\
4 & 000100 & 16 & Yù & 36 & 100100 & 51 & Zhèn \\
5 & 000101 & 35 & Jìn & 37 & 100101 & 21 & Shì Hé \\
6 & 000110 & 45 & Cuì & 38 & 100110 & 17 & Suí \\
7 & 000111 & 12 & Pǐ & 39 & 100111 & 25 & Wú Wàng \\
8 & 001000 & 15 & Qiān & 40 & 101000 & 36 & Míng Yí \\
9 & 001001 & 52 & Gèn & 41 & 101001 & 22 & Bì \\
10 & 001010 & 39 & Jiǎn & 42 & 101010 & 63 & Jì Jì \\
11 & 001011 & 53 & Jiàn & 43 & 101011 & 37 & Jiā Rén \\
12 & 001100 & 62 & Xiǎo Guò & 44 & 101100 & 55 & Fēng \\
13 & 001101 & 56 & Lǚ & 45 & 101101 & 30 & Lí \\
14 & 001110 & 31 & Xián & 46 & 101110 & 49 & Gé \\
15 & 001111 & 33 & Dùn & 47 & 101111 & 13 & Tóng Rén \\
16 & 010000 & 7 & Shī & 48 & 110000 & 19 & Lín \\
17 & 010001 & 4 & Méng & 49 & 110001 & 41 & Sǔn \\
18 & 010010 & 29 & Kǎn & 50 & 110010 & 60 & Jié \\
19 & 010011 & 59 & Huàn & 51 & 110011 & 61 & Zhōng Fú \\
20 & 010100 & 40 & Jiě & 52 & 110100 & 54 & Guī Mèi \\
21 & 010101 & 64 & Wèi Jì & 53 & 110101 & 38 & Kuí \\
22 & 010110 & 47 & Kùn & 54 & 110110 & 58 & Duì \\
23 & 010111 & 6 & Sòng & 55 & 110111 & 10 & Lǚ \\
24 & 011000 & 46 & Shēng & 56 & 111000 & 11 & Tài \\
25 & 011001 & 18 & Gǔ & 57 & 111001 & 26 & Dà Chù \\
26 & 011010 & 48 & Jǐng & 58 & 111010 & 5 & Xū \\
27 & 011011 & 57 & Xùn & 59 & 111011 & 9 & Xiǎo Chù \\
28 & 011100 & 32 & Héng & 60 & 111100 & 34 & Dà Zhuàng \\
29 & 011101 & 50 & Dǐng & 61 & 111101 & 14 & Dà Yǒu \\
30 & 011110 & 28 & Dà Guò & 62 & 111110 & 43 & Guài \\
31 & 011111 & 44 & Gòu & 63 & 111111 & 1 & Qián \\
\bottomrule
\end{tabular}
\end{table}


% ═══════════════════════════════════════════
\section{Cycle Decomposition}\label{sec:cycles}

\subsection{Algorithm}

We apply the standard cycle-finding algorithm: starting from an unvisited element $i$, follow the chain $i \to \sigma(i) \to \sigma^2(i) \to \cdots$ until returning to~$i$, recording the cycle.
Repeat for all 64 elements.

\subsection{Result}

\begin{theorem}\label{thm:main}
The permutation $\sigma$ defined by Equation~\eqref{eq:sigma} has cycle type $\mathbf{(52,\, 10,\, 2)}$ with zero fixed points.
\end{theorem}

\begin{proof}
Direct computation from the data in Table~\ref{tab:mapping}.
Reproducible code is provided in Appendix~A.
\end{proof}

The three cycles are:

\medskip
\noindent\textbf{Primary cycle (length 52, 81.25\%):}
\begin{quote}\small
Qián $\to$ Kūn $\to$ Bō $\to$ Kùn $\to$ Gé $\to$ Lín $\to$ Kǎn $\to$ Héng $\to$ Gòu $\to$ Jiā Rén $\to$ Zhèn $\to$ Jié $\to$ Xiǎo Chù $\to$ Qiān $\to$ Xián $\to$ Dà Guò $\to$ Xùn $\to$ Tài $\to$ Jiǎn $\to$ Suí $\to$ Shī $\to$ Cuì $\to$ Fēng $\to$ Duì $\to$ Dà Chù $\to$ Gǔ $\to$ Méng $\to$ Guān $\to$ Huàn $\to$ Xū $\to$ Yù $\to$ Dùn $\to$ Fù $\to$ Sòng $\to$ Jìn $\to$ Zhūn $\to$ Bǐ $\to$ Pǐ $\to$ Jiàn $\to$ Guī Mèi $\to$ Kuí $\to$ Shì Hé $\to$ Jiě $\to$ Wú Wàng $\to$ Shēng $\to$ Lí $\to$ Dǐng $\to$ Sǔn $\to$ Míng Yí $\to$ Yì $\to$ Bì $\to$ Wèi Jì $\to$ [back to Qián]
\end{quote}

\noindent\textbf{Secondary cycle (length 10, 15.625\%):}
\begin{quote}\small
Lǚ (10) $\to$ Gèn $\to$ Zhōng Fú $\to$ Dà Zhuàng $\to$ Yí $\to$ Jǐng $\to$ Tóng Rén $\to$ Xiǎo Guò $\to$ Dà Yǒu $\to$ Lǚ (56) $\to$ [back to Lǚ (10)]
\end{quote}

\noindent\textbf{Transposition (length 2, 3.125\%):}
\begin{quote}\small
Guài $\leftrightarrow$ Jì Jì
\end{quote}

\subsection{Key Properties}

\begin{table}[ht]
\centering
\caption{Structural properties of $\sigma$.}
\label{tab:properties}
\begin{tabular}{@{}lrl@{}}
\toprule
Property & Value & Interpretation \\
\midrule
Number of cycles & 3 & Few (random expectation $\approx 4.7$) \\
Fixed points & 0 & No hexagram stays in its natural binary position \\
Largest cycle fraction & 81.25\% & High coupling \\
Order of $\sigma$ & $\mathrm{lcm}(52,10,2) = 260$ & $\sigma^{260} = \mathrm{id}$ \\
\bottomrule
\end{tabular}
\end{table}


% ═══════════════════════════════════════════
\section{Statistical Context}\label{sec:stats}

To assess whether the cycle type $(52,10,2)$ is unusual, we compare it with the distribution of cycle structures in random permutations of 64~elements.

\subsection{Fixed Points}

The number of fixed points in a random permutation of $n$ elements is approximately Poisson with $\lambda = 1$.
Thus $P(\text{0 fixed points}) = e^{-1} \approx 0.368$.
Zero fixed points alone is not statistically remarkable.

\subsection{Number of Cycles}

The expected number of cycles in a random permutation of $n$ elements is the $n$-th harmonic number $H_n$.
For $n = 64$, $H_{64} \approx 4.74$, with standard deviation $\approx 2.0$.
The observed value of 3 cycles is at approximately $-0.9\sigma$---mildly low but not extreme on its own.

\subsection{Largest Cycle Length}

The expected length of the longest cycle in a random permutation of $n$ elements is approximately $\lambda n$ where $\lambda \approx 0.6243$ (the Golomb--Dickman constant), giving an expected value of $\approx 40$ for $n=64$.
The observed value of 52 is in the upper tail of the distribution.

\subsection{Monte Carlo Estimate}

A simulation of $2 \times 10^6$ random permutations of 64 elements yields:
\begin{itemize}
  \item Probability of cycle type \emph{exactly} $(52,10,2)$: $\approx 0.094\%$ ($\sim$1 in 1{,}064).
  \item Probability of exactly 3 cycles: $\approx 16.2\%$.
  \item Probability of largest cycle $\geq 52$: $\approx 5.8\%$.
\end{itemize}

\begin{remark}
  The exact cycle type $(52,10,2)$ is uncommon ($\sim$0.1\%) but not extraordinarily rare.
  We refrain from claiming statistical ``significance'' in a hypothesis-testing sense, as there is no well-defined null hypothesis for why an ancient text should follow any particular distribution.
  The value of the decomposition lies in its descriptive power, not in a $p$-value.
\end{remark}

We emphasize that any specific cycle type of a random permutation is individually rare, and we do not claim that 0.094\% constitutes evidence of intentional design. The value of this decomposition is structural and descriptive: it reveals that 81.25\% of hexagrams form a single connected orbit, meaning that iterating $\sigma$ cycles through nearly the entire set before returning. This is analogous to reporting that a molecule forms one large ring rather than several small chains---a structural fact, not a statistical claim.


% ═══════════════════════════════════════════
\section{Hamming Distance Analysis}\label{sec:hamming}

\subsection{Definition}

For two hexagrams encoded as 6-bit strings, the \emph{Hamming distance} $d$ is the number of positions at which the bits differ.

We compute the sequence of Hamming distances $d_1, d_2, \ldots, d_{63}$ between consecutive hexagrams in the King Wen order and compare with the binary natural order and a random baseline.

\subsection{Results}

\begin{table}[ht]
\centering
\caption{Hamming distance statistics for three orderings.}
\label{tab:hamming}
\begin{tabular}{@{}lccc@{}}
\toprule
Statistic & King Wen & Binary natural & Random baseline \\
\midrule
Mean Hamming distance & \textbf{3.349} & 1.905 & 3.000 \\
Pairs with $d=6$ (total inversion) & 9 & 0 & $\sim$1.4 \\
Pairs with $d \leq 2$ (gradual) & 22 & 53 & $\sim$19 \\
Pairs with $d \geq 4$ (contrast) & 28 & 1 & $\sim$21 \\
Even $d$ : odd $d$ ratio & \textbf{3.2 : 1} & --- & 1 : 1 \\
\bottomrule
\end{tabular}
\end{table}

\subsection{Discussion}

Three features stand out:

\begin{enumerate}
  \item \textbf{Contrast bias ($\bar{d} > 3.0$).}
  The mean Hamming distance of 3.349 exceeds the random baseline of 3.0, indicating that the King Wen sequence tends to place structurally dissimilar hexagrams adjacent to each other---a ``contrast narrative'' strategy.

  \item \textbf{Even--odd asymmetry (3.2 : 1).}
  Steve Moore previously noted that even Hamming distances predominate in the King Wen sequence.
  We independently confirm this ratio (76.2\% even vs.\ 23.8\% odd) and identify its mathematical origin: the 32 pairs of \emph{inverted/complementary hexagrams} (\textit{zōng guà}/\textit{cuò guà}) that structure the King Wen sequence naturally produce even Hamming distances (complementary hexagrams always have $d=6$; inverted hexagrams have even~$d$).

  \item \textbf{Nine total inversions ($d=6$).}
  The random expectation is $\sim$1.4 pairs. Nine total inversions means the King Wen sequence frequently places a hexagram next to its bitwise complement---the mathematical expression of the traditional \emph{cuò guà} (complementary hexagram) concept.
\end{enumerate}


% ═══════════════════════════════════════════
\section{Semantic Analysis and Its Limits}\label{sec:semantic}

\subsection{The Ten-Hexagram Cycle}

The secondary cycle contains ten hexagrams which, if their traditional names are read in cycle order, suggest a narrative of personal cultivation:

\begin{quote}\small
Lǚ (treading/propriety) $\to$ Gèn (stillness) $\to$ Zhōng Fú (inner truth) $\to$ Dà Zhuàng (great power) $\to$ Yí (nourishment) $\to$ Jǐng (the well) $\to$ Tóng Rén (fellowship) $\to$ Xiǎo Guò (small exceeding) $\to$ Dà Yǒu (great possession) $\to$ Lǚ (the wanderer)
\end{quote}

\subsection{The Transposition}

Guài (``breakthrough/resolution'') and Jì Jì (``completion'')---one pointing to a decisive beginning, the other to a finished state---swap positions, forming a symbolic loop between initiation and consummation.

\subsection{Methodological Warning}

\begin{remark}[Post-hoc rationalization]\label{rem:posthoc}
  The above semantic readings carry a high risk of \emph{apophenia} (pattern perception in noise).
  The 64 hexagram names are sufficiently evocative that \emph{almost any} permutation's cycles could be narrated as a meaningful sequence.
  This is a cognitive bias, not a mathematical conclusion.

  \textbf{What is verifiable:} the cycle type $(52,10,2)$, zero fixed points, and Hamming distance statistics.

  \textbf{What is not verifiable:} that the ten-hexagram cycle ``means'' personal growth.

  This paper strictly separates the two.
\end{remark}


% ═══════════════════════════════════════════
\section{Related Work}\label{sec:related}

\begin{table}[ht]
\centering
\caption{Prior work on mathematical aspects of the \textit{I Ching} hexagram orderings.}
\label{tab:related}
\begin{tabular}{@{}lll@{}}
\toprule
Author & Contribution & Relation to this paper \\
\midrule
Leibniz (1703) & Binary interpretation of Fu Xi order & Foundational premise \\
Shao Yong (11th c.) & Reconstruction of binary natural order & Foundational premise \\
\citet{moore} & Even:odd Hamming ratio $\approx 3\!:\!1$ & Independently confirmed \\
Cook (2006) & Paired-structure analysis & Different approach (no $S_n$) \\
Petoukhov (2008) & Algebraic structures in hexagrams & Complementary perspective \\
Various & \textit{Zōng guà}/\textit{cuò guà} pair analysis & Reexamined via permutation group \\
\bottomrule
\end{tabular}
\end{table}

To our knowledge, no prior work has performed a cycle decomposition of the permutation mapping between these two orderings.
As of March 2026, searches in Google Scholar (``King Wen sequence'' + ``permutation cycle''), OEIS, CNKI (中国知网: ``文王卦序'' + ``置换''), and major monographs on \textit{I Ching} mathematics have returned no results reporting the cycle type $(52,10,2)$.

This does not prove absolute novelty---only that the result is absent from the publicly searchable literature.

The traditional explanation for the King Wen sequence is found in the \textit{X\`u Gu\`a} (序卦传, Sequence of Hexagrams commentary), one of the Ten Wings. A third historical ordering, found on the Mawangdui silk manuscript (excavated 1973), organizes hexagrams by upper trigram. Applying the same cycle decomposition to other pairs of orderings is a natural extension of this work.


% ═══════════════════════════════════════════
\section{Conclusion}\label{sec:conclusion}

\begin{enumerate}
  \item \textbf{Mathematical fact.} The permutation from the binary natural order to the King Wen sequence has cycle type $(52,10,2)$ with zero fixed points.

  \item \textbf{Structural observation.} The cycle type (52, 10, 2) is uncommon ($\sim$0.1\%) but not extraordinarily rare. The observation of interest is that 81.25\% of elements form a single orbit, indicating high coupling. The mean Hamming distance (3.349 vs.\ baseline 3.0) and even:odd ratio (3.2:1) further characterize the reordering.

  \item \textbf{Hamming distance.} The even:odd ratio of 3.2:1 and the 9 total inversions are consequences of the inverted/complementary pair structure of the King Wen sequence.

  \item \textbf{Epistemological boundary.} Cycle structure is fact; semantic interpretation of cycle membership is conjecture. The two must not be conflated.

  \item \textbf{Historical nuance.} The binary natural order was formalized $\sim$2{,}100 years after the King Wen sequence. Our analysis characterizes the permutation between two independently defined orderings---it does not presuppose that one was derived from the other.
\end{enumerate}


% ═══════════════════════════════════════════
\section*{Acknowledgments}
The author thanks Claude (Anthropic) for assistance with computation and manuscript preparation.
The critical observation that the ``Fu Xi sequence'' framing involves an unverified historical assumption arose during human--AI dialogue and led to the corrected epistemological framework presented in this paper.


% ═══════════════════════════════════════════
\appendix
\section{Reproducible Code}\label{app:code}

The following JavaScript code computes the cycle decomposition from the raw mapping data.
It can be executed in any modern browser console or Node.js environment.

\begin{lstlisting}[caption={Cycle decomposition of $\sigma$.}]
// King Wen number (1-indexed) for each binary natural position (0-indexed)
const W = [
  2,23,8,20,16,35,45,12, 15,52,39,53,62,56,31,33,
  7,4,29,59,40,64,47,6,  46,18,48,57,32,50,28,44,
  24,27,3,42,51,21,17,25, 36,22,63,37,55,30,49,13,
  19,41,60,61,54,38,58,10, 11,26,5,9,34,14,43,1
];

// Build permutation: sigma(i) = W[i] - 1
const sigma = W.map(w => w - 1);

// Cycle decomposition
const visited = new Set(), cycles = [];
for (let i = 0; i < 64; i++) {
  if (visited.has(i)) continue;
  const cycle = [];
  let j = i;
  while (!visited.has(j)) {
    visited.add(j);
    cycle.push(j);
    j = sigma[j];
  }
  cycles.push(cycle);
}
cycles.sort((a, b) => b.length - a.length);

// Output: cycle lengths
console.log(cycles.map(c => c.length));
// => [52, 10, 2]
\end{lstlisting}


\section{Complete Cycle Membership}\label{app:cycles}

\subsection*{Primary Cycle (52 hexagrams)}

{\small
\noindent
Qián (1) $\to$ Kūn (2) $\to$ Bō (23) $\to$ Kùn (47) $\to$ Gé (49) $\to$ Lín (19) $\to$ Kǎn (29) $\to$ Héng (32) $\to$ Gòu (44) $\to$ Jiā Rén (37) $\to$ Zhèn (51) $\to$ Jié (60) $\to$ Xiǎo Chù (9) $\to$ Qiān (15) $\to$ Xián (31) $\to$ Dà Guò (28) $\to$ Xùn (57) $\to$ Tài (11) $\to$ Jiǎn (39) $\to$ Suí (17) $\to$ Shī (7) $\to$ Cuì (45) $\to$ Fēng (55) $\to$ Duì (58) $\to$ Dà Chù (26) $\to$ Gǔ (18) $\to$ Méng (4) $\to$ Guān (20) $\to$ Huàn (59) $\to$ Xū (5) $\to$ Yù (16) $\to$ Dùn (33) $\to$ Fù (24) $\to$ Sòng (6) $\to$ Jìn (35) $\to$ Zhūn (3) $\to$ Bǐ (8) $\to$ Pǐ (12) $\to$ Jiàn (53) $\to$ Guī Mèi (54) $\to$ Kuí (38) $\to$ Shì Hé (21) $\to$ Jiě (40) $\to$ Wú Wàng (25) $\to$ Shēng (46) $\to$ Lí (30) $\to$ Dǐng (50) $\to$ Sǔn (41) $\to$ Míng Yí (36) $\to$ Yì (42) $\to$ Bì (22) $\to$ Wèi Jì (64) $\to$ [Qián]
}

\subsection*{Secondary Cycle (10 hexagrams)}

{\small
\noindent
Lǚ-10 (10) $\to$ Gèn (52) $\to$ Zhōng Fú (61) $\to$ Dà Zhuàng (34) $\to$ Yí (27) $\to$ Jǐng (48) $\to$ Tóng Rén (13) $\to$ Xiǎo Guò (62) $\to$ Dà Yǒu (14) $\to$ Lǚ-56 (56) $\to$ [Lǚ-10]
}

\subsection*{Transposition}

{\small
\noindent
Guài (43) $\leftrightarrow$ Jì Jì (63)
}


% ═══════════════════════════════════════════
\begin{thebibliography}{9}

\bibitem[Cook(2006)]{cook2006}
Cook, R.~S.
\textit{Classical Chinese Combinatorics: Derivation of the Book of Changes Hexagram Sequence}.
Ph.D.\ dissertation, University of California, Berkeley, 2003. Published as STEDT Monograph 5, 2006.

\bibitem[Leibniz(1703)]{leibniz1703}
Leibniz, G.~W.
Explication de l'Arithm{\'e}tique Binaire.
\textit{M{\'e}moires de l'Acad{\'e}mie Royale des Sciences}, 1703.

\bibitem[Petoukhov(2008)]{petoukhov2008}
Petoukhov, S.~V.
\textit{Matrix Genetics, Algebras of the Genetic Code, Noise-Immunity}.
RCD, 2008.

\bibitem[Rutt(2002)]{rutt2002}
Rutt, R.
\textit{The Book of Changes (Zhouyi): A Bronze Age Document}.
Routledge, 2002.

\bibitem[Shaughnessy(2014)]{shaughnessy2014}
Shaughnessy, E.~L.
\textit{Unearthing the Changes: Recently Discovered Manuscripts of the Yi Jing}.
Columbia University Press, 2014.

\bibitem[Moore]{moore}
Moore, S.
Structural Elements in the King Wen Sequence of Hexagrams.
Unpublished manuscript, available at \url{http://biroco.com/yijing/}.

\end{thebibliography}

\end{document}
